\section{Introduction}\label{introduction}

In the model presented here, the phenomenon of gravity is not introduced
as a fundamental force, but rather emerges from the internal dynamics of
a tensioned medium --- a three-dimensional brane embedded in a
higher-dimensional space. This use of the term ``brane'' should not be confused with Randall--Sundrum--type brane-world scenarios, in which a thin 3+1-dimensional brane with tuned tension sits in a warped higher-dimensional Einstein-gravity bulk; in the present work, the brane is a physical elastic medium embedded in flat four-dimensional space, and gravity is interpreted as emergent from its induced metric and internal tension dynamics. A key feature of this model is that
deformations in the brane's fourth (amplitude) dimension are not
independent of its spatial structure: instead, the brane resists local
stretching by contracting in the lateral spatial directions. This
\textbf{geometric coupling} between amplitude fluctuations and lateral
tension causes the brane to adjust its shape in such a way that energy
concentrations --- represented by oscillations or solitonic structures
in the amplitude dimension --- induce curvature-like behavior in the
three spatial dimensions. The resulting contraction fields resemble the
attractive behavior attributed to gravity in general relativity, yet
they arise here from local tension dynamics within a medium governed by
deterministic laws. Thus, gravity appears as a \textbf{secondary,
emergent effect} of the system's effort to minimize internal deformation
--- a direct consequence of the coupling between different embedding
directions of the brane.

A central heuristic in this work is that \emph{mass is nothing but
localized energy in motion}, in line with the standard relativistic
equivalence between inertial and gravitational mass. In particular, the
gravitational mass of any closed system is taken to be proportional to
its total energy in the center-of-mass frame, including kinetic and
internal degrees of freedom. This viewpoint is emphasized clearly by
van der Mark and 't~Hooft in their essay \emph{``Light is Heavy''},
where they show that a box filled with radiation weighs more than an
empty one, and that the gravitational mass of the light inside is
precisely $E/c^2$ even though each photon has zero rest mass
\cite{vanderMark2000LightHeavy}. The present brane model simply
implements this idea mechanically: all mass, whether ``matter'' or
``radiation'', is reinterpreted as internal energy stored in waves and
tension of a single medium.

The present work builds on the toroidal electron model of Williamson and van der Mark (1997), generalizing it into a continuous brane framework where confinement, relativity, and gravitation emerge from a unified geometric substrate (detailed in Section~\ref{internal-structure-of-the-electron}).

\subsection{Core Idea and Scope}\label{core-idea}

The central hypothesis explored here is that \textbf{lateral contraction} in a tensioned medium—typically neglected under the small-angle approximation in classical wave theory—may, when extended to higher dimensions, provide a natural mechanism for gravitational effects. Specifically, if a 3D brane is embedded in 4D space with coupling between amplitude (fourth-direction) deformations and in-brane tension, localized energy concentrations would induce spatially extended contraction patterns resembling gravitational fields.

This paper develops that idea into a coherent framework in which:
\begin{itemize}
\item \textbf{Special relativity} emerges from isotropic wave propagation in an absolute medium;
\item \textbf{Quantum uncertainty} reflects Fourier-localization constraints on wave packets;
\item \textbf{Particles} correspond to topologically stable solitons with internal phase structure;
\item \textbf{Electromagnetism} arises from Lorentz transformations of a scalar interaction;
\item \textbf{Gravitation} emerges from amplitude-induced curvature of the brane's induced metric.
\end{itemize}

While speculative and exploratory, the framework aims to demonstrate that a single elastic substrate can, in principle, unify phenomena typically treated as separate in standard physics.

\paragraph{Microscopic time vs emergent spacetime.}
Throughout this work we distinguish between a microscopic evolution parameter $t$ and the emergent spacetime experienced by observers. At the fundamental level, the brane configuration
\[
\mathbf X : \Omega \times \mathbb{R} \to \mathbb{R}^4, \qquad (x,t) \mapsto \mathbf X(x,t)
\]
evolves with respect to an external time parameter $t$ that is not part of the embedding itself. At the emergent level, observers reconstruct effective spacetime coordinates $x^\mu = (ct, x^i)$ and an approximate Lorentzian metric $g_{\mu\nu}(x)$ from the propagation of waves on the brane. In other words, relativity arises as an effective description of wave kinematics on a pre-existing brane evolving in microscopic time $t$.

\paragraph{Structure of this paper.}
The structure of the paper mirrors this substrate–first perspective.
Section~\ref{conceptual-model} defines the microscopic brane substrate, its
embedding, Lagrangian, and continuum equations of motion.
Section~\ref{reconstructing-physics-as-an-emergent-brane-phenomenon} then treats
familiar physical theories — special relativity, quantum wave mechanics,
electromagnetism, gravity, particle phenomenology, and measurement — as emergent,
effective descriptions reconstructed from that substrate.
Section~\ref{experimental-setting} describes a discrete simulation framework that
integrates only the brane degrees of freedom while reconstructing emergent fields
and observables as measured quantities.
Finally, Sections~\ref{discussion} and~\ref{conclusion} discuss open problems,
limitations, and the broader implications of this reconstruction programme.

\subsection{Related Works}\label{related-works}

Recent work by \textbf{Roth et al.}~\cite{Roth2023QQM2} and \textbf{Close}~\cite{Close2025EquationEverything}
develops closely related ``elastic medium'' approaches. Danielewski, Sapa and Roth
introduce a \emph{quaternion-based elastic-aether model} in which the vacuum is treated
as an ideal elastic solid, particles appear as standing-wave solitons, and gravity
emerges from tension-induced changes in the effective refractive index of the medium,
leading to Dirac-like dynamics from the underlying continuum~\cite{Roth2023QQM2}.
Similarly, Close analyzes spin–angular-momentum density waves in an ideal elastic
solid and derives a Dirac bispinor description as an effective first-order wave
equation for these excitations~\cite{Close2025EquationEverything}. These works
dovetail with the present brane-tension picture, in which relativistic matter fields
and gravitational effects arise from the kinematics of a single oscillatory medium.

More broadly, our interpretation fits into the established programme of
\emph{emergent gravity} and \emph{analogue spacetime} models. Sakharov's classic
proposal treats the Einstein--Hilbert action as an induced elastic response of the
quantum vacuum~\cite{Sakharov1967VacuumGravity}. Our three-dimensional
brane can be viewed as a specific realization of this general idea: a tensioned
medium whose low-energy excitations admit an effective relativistic field theory
description.

Finally, there exist modern ``aether-like'' theories, such as Einstein--aether
models, which introduce a dynamical preferred-frame field coupled to gravity and
study systematic violations of local Lorentz symmetry within a generally covariant
framework~\cite{Jacobson2004EinsteinAether}. While our brane model posits an
elastic medium rather than a vector aether field, it shares with these approaches
the basic premise that relativistic spacetime symmetries may be emergent rather
than fundamental.

A complementary line of work studies the internal structure of electrons and
photons in terms of structured electromagnetic waves on Compton scales. The
toroidal model of Williamson and van der Mark treats the electron as a bound
photon circulating on a torus of radius $\mathcal{O}(\lambda_C)$ and uses
internal energy density and geometry to estimate the observed charge
magnitude~\cite{Williamson1997,Williamson2015}. More recently, zitterbewegung
approaches place the characteristic radius of the internal trembling motion at
roughly half the reduced Compton wavelength, again highlighting the special role
of the Compton scale in electron structure~\cite{UrdanetaSantos2023}. On the
quantum-optics side, detailed analyses of single-photon wave equations and field
normalisation~\cite{KiesslingTahvildarZadeh2018SinglePhoton,Liu2018SinglePhotonFields}
make explicit how the classical wave energy of a mode is tied to the quantum
energy~$\hbar\omega$. Both strands inform our use of Compton-scale solitons and
our procedure for fixing the amplitude normalisation of brane waves from the
condition $E_{\text{wave}}=\hbar\omega$ for a single quantum.